\section{实战演练：复现讲座精髓}

\subsection{素材与准备}
复现这节讲座需要先拿到相关资产：视频 `audio.m4a`、`lecture09.srt`、封面 `cover.jpg` 以及任何可用的 slides/handout。把字幕清洗成 \texttt{subs\_clean.txt}（去掉空行、重复句）后，就可以根据时间戳快速定位讲师的重点段落。Metadata 也要准备完整：作者、发布日期、视频链接与时长，方便前置模板在封面箱里直接引用。若 slides 已经披露，可以用 `magick -density 150 slide.pdf slide.jpg` 把每页转换成图，放在 `figure` 中补充视觉证据；若没有 slides，就必须从视频截取关键帧，并写明具体 timecode。

\begin{knowledgebox}{Slides 与关键帧的准备}
\begin{itemize}
    \item 若 slides 已有 PDF，优先截取重要页并插入 `figure`，同时在 caption 里注明页码与时间指示。
    \item 如果没有 slides，就在视频里找到关键 moment（如 00:22:10 或 00:35:40）截帧，放入 `figure` 并注释时间。
    \item 标注 `\textit{}` 原话时将时间戳写成 `[HH:MM:SS]` 形式，便于 later reference。
\end{itemize}
\end{knowledgebox}

\subsection{拆解与写作流程}
写笔记的流程可以分为：读取 clean subtitles、梳理 ``Teaching Logic Outline''、按照逻辑扩展每一节、插入 boxes、加上 figure 与 summary。每个 section 最终都配 `\subsection{本章小结}`，形成可复现的 structure。
\begin{itemize}
    \item 自 \texttt{subs\_clean} 提取主题关键词，按 Data / Algorithm / Agent / Infra 等主题排序。
    \item 用 `\textit{''}` 引用讲师原话，标注时间戳并写入对应小节。
    \item 把关键洞察写成 `importantbox`、`knowledgebox`、`warningbox`，保证不同类型的信息有清晰视觉层次。
    \item 最后用 summary table + further reading 摘要整期讲稿，便于复现者快速复盘。
\end{itemize}
\begin{knowledgebox}{复现笔记时的写作提示}
\begin{itemize}
    \item 以教学逻辑组织章节，不要只照字幕的时间顺序堆内容。
    \item 每一个重要观点都要复述、加翻译并标注出处，用 `\textit{}` 保留原味。
    \item Visual evidence 要附 figure + time footnote，cover、slides、关键帧都可以。
    \item `warningbox` 用来提醒潜在误读，`importantbox` 用来强调突破性结论。
\end{itemize}
\end{knowledgebox}

\subsection{视觉证据与合成}
本节没有 slides，所以我们只能用 video frame 作为图片来源：封面 `cover.jpg`、speaker-side 演示、工具调用界面。每张 figure 附上 `\protect\footnotemark` 说明时间区间（例如 00:00:05--00:00:13），并在正文里解释这张图为什么关键。若未来找到 slides，可以用 `magick -density ...` 把 PDF 转成 `slide-*.jpg`，再把它们插入到 `figure` 中替代 frame。每个 timecode 都要在 caption 里写清楚，以便后续 audit 或 review 时验证。

\begin{importantbox}{视觉证据的高质量要求}
撷取关键帧图时务必保持 150dpi 以上的清晰度，caption 里注明时间点，并对照讲师原话解释画面的意义。若用 slides，记得对齐页码与时间线，避免留下未替换的链接占位文本或空白 caption。
\end{importantbox}

\subsection{复现流程与关键帧}
复现流程可以分成三个阶段：1) 从 \texttt{subs\_clean} 里截取主题，2) 按照 Teaching Logic 把主题分配到各个 section，3) 插入图文/box/summary。每个 section 都应记录对应的 timestamp，类似 `[00:40:12]` 对应 Agent Swarm 的 critical steps 说明。若发现视频里有 slides 但字幕没提到，仍可手动对其翻译，写入 `figure` 说明里，并注明 ``Slide 暂缺英文解释''。

\begin{knowledgebox}{关键帧与 audit-ready 文档}
\begin{itemize}
    \item 每个 figure 后面都写上 timecode，例如 ``画面时间区间：00:32:10--00:32:48''。
    \item 摘要表格里标明 topics 与 evidence，如 ``Agent Swarm -> 4.5 倍 efficiency -> [00:36:50] figure''。
    \item 避免使用未替换的 metadata 占位文本，所有封面信息都必须写实。
\end{itemize}
\end{knowledgebox}

\subsection{本章小结}
这套实战流程：先收集素材并整理 metadata，再按教学逻辑拆解，再对每章补充视觉证据、引用与 boxes，就能产出符合质量标准的 K2.5 笔记。

